\documentclass[11pt,a4paper]{article}

% -------------------- Langue / encodage --------------------
\usepackage[T1]{fontenc}
\usepackage[french]{babel}
\usepackage{array}

% -------------------- Mise en page --------------------
\usepackage{geometry}
\geometry{
  a4paper,
  landscape,
  left=1.6cm,
  right=1.6cm,
  top=1.0cm,
  bottom=1.2cm
}


% -------------------- Maths / mise en forme --------------------
\usepackage{amsmath,amssymb}
\usepackage{parskip}
\usepackage[explicit]{titlesec}
\usepackage[normalem]{ulem} % pour \uline sans casser \emph
\usepackage{titlesec}

\usepackage{fancyhdr}
\usepackage{lastpage}
\usepackage[hidelinks]{hyperref}
\usepackage{enumitem}
\usepackage[table]{xcolor} % <-- IMPORTANT : pour \rowcolor dans les tableaux
\usepackage{paracol}

% -------------------- TikZ + tkz-tab --------------------
\usepackage{tikz}
\usetikzlibrary{arrows.meta,calc}
\usepackage{tkz-tab}

% -------------------- Variables --------------------
\newcommand{\FicheTitre}{Logarithme népérien -- Fiche de cours}
\newcommand{\Niveau}{Terminale générale -- Mathématiques}
\newcommand{\Annee}{Année scolaire 2025/2026}

\newcommand{\SiteURL}{https://naoufel-coursparticuliers.netlify.app/}
\newcommand{\SiteTexte}{naoufel-coursparticuliers.netlify.app}

% -------------------- Footer --------------------
\pagestyle{fancy}
\fancyhf{}
\lfoot{\textit{\Niveau \;--\; \Annee}}
\cfoot{\thepage/\pageref{LastPage}}
\rfoot{\href{\SiteURL}{\texttt{\SiteTexte}}}
\renewcommand{\headrulewidth}{0pt}
\renewcommand{\footrulewidth}{0.4pt}

% -------------------- Spacing (anti débordement) --------------------
\setlength{\parskip}{4pt}
\setlength{\parindent}{0pt}
\setlist[itemize]{leftmargin=*, itemsep=2pt, topsep=2pt}


% -------------------- Titres --------------------
%\titleformat{\section}{\large\bfseries}{\thesection.}{0.5em}{} %noir
%\titleformat{\subsection}{\normalsize\bfseries}{\alph{subsection}.}{0.5em}{}\titleformat{\section}
% --- TITRES colorés + petite ligne sous le TEXTE (sans souligner le numéro) ---
% SECTION : numéro rouge (non souligné) + titre rouge souligné
\titleformat{\section}[hang]
  {\normalsize\bfseries}
  {\textcolor{red}{\thesection.}}
  {0.2em}
  {\textcolor{red}{\uline{#1}}}

% SUBSECTION : lettre verte (non soulignée) + titre vert souligné
\titleformat{\subsection}[hang]
  {\normalsize\bfseries}
  {\hspace*{1.2em}\textcolor{green!60!black}{\thesubsection}}
  {0.4em}
  {\textcolor{green!60!black}{\uline{#1}}}
% -------------------- Bandeau titre --------------------
\newcommand{\TitreBandeau}[1]{%
  \noindent\hrule
  \vspace{0.55em}
  \begin{center}
    {\LARGE\bfseries #1}
  \end{center}
  \vspace{0.55em}
  \noindent\hrule
  \vspace{0.9em}
}

% -------------------- Helpers --------------------
\newcommand{\defi}{\textbf{Définition :} }
\newcommand{\prop}{\textbf{Propriété :} }
\newcommand{\rem}{\textbf{Remarque :} }

% ============================
% Graphique "grand" de y = ln(x)
% ============================
\newcommand{\GraphLn}{%
\begin{center}
\resizebox{0.65\linewidth}{!}{% <-- plus grand que 0.95
\begin{tikzpicture}[>=Latex, every node/.style={font=\small}]
  % Fenêtre (un peu plus large/haute)
  \def\xmin{-1} \def\xmax{7}
  \def\ymin{-3.5} \def\ymax{3.5}

  % Grille
  \draw[step=1, very thin, blue!25] (\xmin,\ymin) grid (\xmax,\ymax);

  % Axes
  \draw[very thick, ->] (\xmin,0) -- (\xmax+0.2,0);
  \draw[very thick, ->] (0,\ymin) -- (0,\ymax+0.2);
  \node[below left] at (0,0) {$O$};

  % Asymptote x=0
  \draw[dashed, thick] (0,\ymin) -- (0,\ymax);
  \node[above right] at (0,\ymax) {\small $y=ln(x)$};
   % Asymptote y=0
  \draw[dashed, thick] (\xmin,0) -- (\xmax,0);
  \node[above right] at (\xmax,0) {\small $x$};
  

  % Courbe y=ln(x)
  \draw[thick, domain=0.06:\xmax, samples=250]
    plot (\x, {ln(\x)});

  % Points repères : 1, e, 2
  \fill (1,0) circle (2.6pt);
  \node[below] at (1,0) {$1$};


  \fill ({exp(1)},1) circle (2.6pt);
  \node[above right] at ({exp(1)},1) {\small $(e,1)$};

  \node[above right] at (5.2,{ln(5.2)}) {\small $y=\ln(x)$};
\end{tikzpicture}%
}%
\end{center}
}

\begin{document}

\TitreBandeau{\FicheTitre}

% -------------------- 2 colonnes + ligne verticale --------------------
\columnratio{0.5,0.5}
\setlength{\columnsep}{1.2cm}
\setlength{\columnseprule}{0.6pt}
\def\columnseprulecolor{\color{black}}

\begin{paracol}{2}

% ====== COLONNE GAUCHE ======
\section{Définition et propriétés}

\subsection{Définition (fonction réciproque de exp)}
\defi La fonction \textbf{logarithme népérien}, notée $\ln$, est définie sur
$\mathbb{R}_+^*=]0,+\infty[$ et vérifie :
\[
\ln(x)=y \iff e^{\,y}=x
\]
\subsection{Propriétés de calcul}
\prop Pour $a>0$ et $b>0$ :
\begin{itemize}[nosep]
  \item $\ln(ab)=\ln(a)+\ln(b)$
  \item $\ln\!\left(\dfrac{a}{b}\right)=\ln(a)-\ln(b)$
  \item $\ln(a^n)=n\ln(a)$ (pour $n\in\mathbb{Z}$)
\end{itemize}

\prop Valeurs à connaître : $\ln(1)=0$ et $\ln(e)=1$.


\subsection{Représentation graphique}
\GraphLn

% >>> Basculer sur la colonne de droite <<<
\switchcolumn
\setcounter{section}{1} % <-- pour que "Variations" devienne 2.

% ====== COLONNE DROITE ======
\section{Variations, signe et limites}

\subsection{Dérivées}
\prop Sur $]0,+\infty[$ :
\[
(\ln x)'=\frac{1}{x}
\]
$\dfrac{1}{x}>0$ sur $]0,+\infty[$, ainsi $\ln$ est \textbf{strictement croissante}.

\prop Si $u(x)>0$ :
\[
(\ln(u(x)))'=\frac{u'(x)}{u(x)}
\]

\subsection{Tableau de variations}
\begin{center}
\begin{tikzpicture}
  \tkzTabInit{$x$/1, $\ln(x)$/1}{$0^{+}$, $1$, $+\infty$}
  \tkzTabVar{-/$-\infty$, R/ $0$, +/$+\infty$}
\end{tikzpicture}
\end{center}

\subsection{Signe de $\ln$}
\begin{center}
\begin{tikzpicture}
  \tkzTabInit{$x$/1, $\ln(x)$/1}{$0^{+}$, $1$, $+\infty$}
  \tkzTabLine{, -, z, +, }
\end{tikzpicture}
\end{center}

\subsection{Limites}
\prop
\[
\lim_{x\to 0^+}\ln(x)=-\infty
\qquad\text{et}\qquad
\lim_{x\to +\infty}\ln(x)=+\infty
\]
\rem Asymptote verticale : $x=0$.

% >>> Basculer sur la colonne de droite <<<
\switchcolumn
\setcounter{subsection}{4} % <-- pour que "CC" devienne e.

\subsection{Croissance comparée}

\prop Pour tout $\alpha>0$ :
\[
\lim_{x\to 0^+} x^{\alpha}\ln(x)=0
\qquad \text{et} \qquad
\lim_{x\to +\infty}\frac{\ln(x)}{x^{\alpha}}=0
\]

\setcounter{section}{2} % <-- pour que "CC" devienne e.

\section{Équations/inéquations avec $\ln$}

\subsection{Méthodes}
\prop Pour résoudre (avec condition $x>0$) :
\[
\ln(x)=k \iff x=e^k
\]
\prop Pour comparer (car $\ln$ est croissante sur $]0,+\infty[$) :
\[
\ln(x)\le \ln(a) \iff x\le a \quad (x>0,\ a>0)
\]


% ====== PAGE SUIVANTE : COLONNE GAUCHE ======
\section{À retenir}
\begin{itemize}
  \item Domaine : $x>0$. Asymptote verticale : $x=0$
  \item $\ln(ab)=\ln a+\ln b$ | $\ln(a/b)=\ln a-\ln b$ | $\ln(a^n)=n\ln a$
  \item Dérivées : $(\ln x)'=1/x$ ; $(\ln(u))'=u'/u$ (si $u>0$)
  \item Signe : $\ln(x)<0$ sur $]0,1[$ ; $\ln(1)=0$ ; $\ln(x)>0$ si $x>1$
  \item Les 2 limites de ln et les 2 croissances comparées
\end{itemize}

\end{paracol}
\end{document}
